% GENERATED FILE. Edit canonical lessons, then run npm run generate:books.
% canonical-source-hash: fnv1a64:36a347c0c90ab361
% canonical-lessons: MW-C214-pherno, MW-C214-sarikho, MW-C214-thaili, MW-C214-palatno, MW-C214-ladhno

\chapter{Taking It Back}
\label{ch:mw-taking-it-back}
\hlchaptermodality{214}

\begin{chapteropening}
\textbf{By the end of this chapter:} \emph{I can take back what I bought and find my receipt.}

Complete the last lesson of chapter 214: I can take back what I bought and find my receipt.
\end{chapteropening}

\section[\texorpdfstring{\mw{फेरणो} (pherṇo)}{फेरणो (pherṇo)}]{\mw{फेरणो} (pherṇo) — to give back; to return (a thing)}

\label{lesson:MW-C214-pherno}



\begin{quote}
Before the new one: say the Marwadi for free of charge, then the Marwadi for to be sold.
\end{quote}

\subsection*{You'll want to know: \mw{फेरणो}}
\textbf{\mw{फेरणो}} — \emph{pherṇo} — \textquotedblleft{}to give back; to return (a thing)\textquotedblright{}.

\begin{quote}
\mw{हूं} \mw{काल} \mw{ये} \mw{कुरतो} \mw{फेरसूं।}
\end{quote}

\begin{quote}
\emph{hū̃ kāl ye kurto phersū̃.}
\end{quote}

\begin{quote}
I'll take this kurta back tomorrow.
\end{quote}

Past (I …): \mw{म्हैं} \mw{कुरतो} \mw{फेर्यो}, \emph{mhaĩ kurto pheryo}.

\begin{cousinweb}
\textbf{\mw{फेरणो}} (\emph{pherṇo}) is inherited from a reconstructed \emph{*phirati}, \textquotedblleft{}turns\textquotedblright{}; CDIAL lists Old Marwari \emph{phiraï}, \textquotedblleft{}turns, walks about\textquotedblright{} (CDIAL 9078). Giving a purchase back is, probably, turning it back; \textbf{\mw{फेर}}, \textquotedblleft{}again\textquotedblright{}, comes from the same word.
\end{cousinweb}

\begin{grammarlens}[title={What you've built}]
One more word for this chapter.
\end{grammarlens}

\subsection*{Your turn}
\begin{itemize}
\raggedright
  \item \emph{Say it:} \emph{pherṇo}
  \item \emph{Say it:} \emph{pherṇo}, once more
  \item \emph{Say it:} \emph{hū̃ kāl ye kurto phersū̃.} — I'll take this kurta back tomorrow.
  \item \emph{Recall:} say the Marwadi for free of charge, then the Marwadi for to be sold, then say \emph{pherṇo} again
\end{itemize}

\begin{tcolorbox}[breakable,colback=teal!4,colframe=teal!35!black,title={Before you move on}]
What does \mw{फेरणो} mean? (\textquotedblleft{}To give back; to return (a thing)\textquotedblright{}.) Say it once more.
\end{tcolorbox}
\section[\texorpdfstring{\mw{सारीखो} (sārīkho)}{सारीखो (sārīkho)}]{\mw{सारीखो} (sārīkho) — alike; the same as}

\label{lesson:MW-C214-sarikho}



\begin{quote}
Before the new one: say the Marwadi for to be sold, then the Marwadi for to give back.
\end{quote}

\subsection*{You'll want to know: \mw{सारीखो}}
\textbf{\mw{सारीखो}} — \emph{sārīkho} — \textquotedblleft{}alike; the same as\textquotedblright{}.

\begin{quote}
\mw{ये} \mw{दो} \mw{कुरता} \mw{सारीखा} \mw{है।}
\end{quote}

\begin{quote}
\emph{ye do kurtā sārīkhā hai.}
\end{quote}

\begin{quote}
These two kurtas are alike.
\end{quote}

\begin{cousinweb}
\textbf{\mw{सारीखो}} (\emph{sārīkho}) is inherited from Sanskrit \emph{sadṛkṣa}, \textquotedblleft{}like\textquotedblright{}; CDIAL lists Old Marwari \emph{sārikho}, \textquotedblleft{}like\textquotedblright{} (CDIAL 13119).
\end{cousinweb}

\begin{grammarlens}[title={What you've built}]
One more word for this chapter.
\end{grammarlens}

\subsection*{Your turn}
\begin{itemize}
\raggedright
  \item \emph{Say it:} \emph{sārīkho}
  \item \emph{Say it:} \emph{sārīkho}, once more
  \item \emph{Say it:} \emph{ye do kurtā sārīkhā hai.} — These two kurtas are alike.
  \item \emph{Recall:} say the Marwadi for to be sold, then the Marwadi for to give back, then say \emph{sārīkho} again
\end{itemize}

\begin{tcolorbox}[breakable,colback=teal!4,colframe=teal!35!black,title={Before you move on}]
What does \mw{सारीखो} mean? (\textquotedblleft{}Alike; the same as\textquotedblright{}.) Say it once more.
\end{tcolorbox}
\section[\texorpdfstring{\mw{थैली} (thailī)}{थैली (thailī)}]{\mw{थैली} (thailī) — a small bag}

\label{lesson:MW-C214-thaili}



\begin{quote}
Before the new one: say the Marwadi for to give back, then the Marwadi for alike.
\end{quote}

\subsection*{You'll want to know: \mw{थैली}}
\textbf{\mw{थैली}} — \emph{thailī} — \textquotedblleft{}a small bag\textquotedblright{}.

\begin{quote}
\mw{सामान} \mw{थैली} \mw{मांय} \mw{राखो।}
\end{quote}

\begin{quote}
\emph{sāmān thailī mā̃y rākho.}
\end{quote}

\begin{quote}
Put the things in the bag.
\end{quote}

A feminine noun.

\begin{cousinweb}
\textbf{\mw{थैली}} (\emph{thailī}) is inherited from Sanskrit \emph{sthavi}, \textquotedblleft{}sack, bag\textquotedblright{}. CDIAL gives no Marwari form in this entry, but lists Hindi \emph{thailī}, \textquotedblleft{}bag\textquotedblright{} (CDIAL 13746).
\end{cousinweb}

\begin{grammarlens}[title={What you've built}]
One more word for this chapter.
\end{grammarlens}

\subsection*{Your turn}
\begin{itemize}
\raggedright
  \item \emph{Say it:} \emph{thailī}
  \item \emph{Say it:} \emph{thailī}, once more
  \item \emph{Say it:} \emph{sāmān thailī mā̃y rākho.} — Put the things in the bag.
  \item \emph{Recall:} say the Marwadi for to give back, then the Marwadi for alike, then say \emph{thailī} again
\end{itemize}

\begin{tcolorbox}[breakable,colback=teal!4,colframe=teal!35!black,title={Before you move on}]
What does \mw{थैली} mean? (\textquotedblleft{}A small bag\textquotedblright{}.) Say it once more.
\end{tcolorbox}
\section[\texorpdfstring{\mw{पलटणो} (palaṭṇo)}{पलटणो (palaṭṇo)}]{\mw{पलटणो} (palaṭṇo) — to flip over; to change}

\label{lesson:MW-C214-palatno}



\begin{quote}
Before the new one: say the Marwadi for alike, then the Marwadi for a small bag.
\end{quote}

\subsection*{You'll want to know: \mw{पलटणो}}
\textbf{\mw{पलटणो}} — \emph{palaṭṇo} — \textquotedblleft{}to flip over; to change\textquotedblright{}.

\begin{quote}
\mw{आप} \mw{थैली} \mw{पलटो।}
\end{quote}

\begin{quote}
\emph{āp thailī palṭo.}
\end{quote}

\begin{quote}
Please turn the bag over.
\end{quote}

Past (I …): \mw{म्हैं} \mw{थैली} \mw{पलटी}, \emph{mhaĩ thailī palṭī}.

\begin{cousinweb}
\textbf{\mw{पलटणो}} (\emph{palaṭṇo}) is inherited from a reconstructed \emph{*pallaṭṭ}, \textquotedblleft{}turn, overturn\textquotedblright{}, whose origin CDIAL calls very doubtful; CDIAL lists Old Marwari \emph{palaṭaï}, \textquotedblleft{}alters\textquotedblright{}, and Hindi \emph{palaṭnā}, \textquotedblleft{}to overturn\textquotedblright{} (CDIAL 7968).
\end{cousinweb}

\begin{grammarlens}[title={What you've built}]
One more word for this chapter.
\end{grammarlens}

\subsection*{Your turn}
\begin{itemize}
\raggedright
  \item \emph{Say it:} \emph{palaṭṇo}
  \item \emph{Say it:} \emph{palaṭṇo}, once more
  \item \emph{Say it:} \emph{āp thailī palṭo.} — Please turn the bag over.
  \item \emph{Recall:} say the Marwadi for alike, then the Marwadi for a small bag, then say \emph{palaṭṇo} again
\end{itemize}

\begin{tcolorbox}[breakable,colback=teal!4,colframe=teal!35!black,title={Before you move on}]
What does \mw{पलटणो} mean? (\textquotedblleft{}To flip over; to change\textquotedblright{}.) Say it once more.
\end{tcolorbox}
\section[\texorpdfstring{\mw{लाधणो} (lādhṇo)}{लाधणो (lādhṇo)}]{\mw{लाधणो} (lādhṇo) — to be found; to turn up}

\label{lesson:MW-C214-ladhno}



\begin{quote}
Before the new one: say the Marwadi for a small bag, then the Marwadi for to flip over.
\end{quote}

\subsection*{You'll want to know: \mw{लाधणो}}
\textbf{\mw{लाधणो}} — \emph{lādhṇo} — \textquotedblleft{}to be found; to turn up\textquotedblright{}.

\begin{quote}
\mw{म्हारी} \mw{रसीद} \mw{कोनी} \mw{लाधै।}
\end{quote}

\begin{quote}
\emph{mhārī rasīd konī lādhai.}
\end{quote}

\begin{quote}
I can't find my receipt.
\end{quote}

Past (it …): \mw{वो} \mw{सिक्को} \mw{लाध्यो}, \emph{vo sikko lādhyo}.

\begin{cousinweb}
\textbf{\mw{लाधणो}} (\emph{lādhṇo}) is inherited from Sanskrit \emph{labdha}, \textquotedblleft{}taken, seized\textquotedblright{}; CDIAL lists the Old Marwari verb \emph{lādhaï}, \textquotedblleft{}finds\textquotedblright{} (CDIAL 10946).
\end{cousinweb}

\begin{grammarlens}[title={What you've built}]
One more word for this chapter.
\end{grammarlens}

\subsection*{Your turn}
\begin{itemize}
\raggedright
  \item \emph{Say it:} \emph{lādhṇo}
  \item \emph{Say it:} \emph{lādhṇo}, once more
  \item \emph{Say it:} \emph{mhārī rasīd konī lādhai.} — I can't find my receipt.
  \item \emph{Recall:} say the Marwadi for a small bag, then the Marwadi for to flip over, then say \emph{lādhṇo} again
\end{itemize}

\begin{tcolorbox}[breakable,colback=teal!4,colframe=teal!35!black,title={Before you move on}]
What does \mw{लाधणो} mean? (\textquotedblleft{}To be found; to turn up\textquotedblright{}.) Say it once more.
\end{tcolorbox}
